\clearpage
\phantomsection% 
\addcontentsline{toc}{chapter}{RINGKASAN}
\thispagestyle{fancy}

\begin{center}
	\large \bfseries \MakeUppercase{Ringkasan}\\
	\normalsize \normalfont {\thetitle}\\
	\normalsize \normalfont {\theauthor}\\
	\bigskip
\end{center}

	\normalsize \normalfont \justifying \singlespacing
	Perkembangan teknologi \textit{Automatic Speech Recognition} (ASR) saat ini telah mencapai kemajuan yang signifikan. Namun, mayoritas sistem ASR yang ada masih berfokus pada bahasa mayor dan kurang mendukung bahasa daerah yang termasuk dalam kategori bahasa dengan sumber daya rendah (\textit{low-resource language}), seperti bahasa Minangkabau. Ketiadaan dataset berskala besar yang teranotasi dengan baik dan kompleksitas variasi dialek menjadikan pemodelan ASR untuk bahasa Minangkabau sebagai sebuah tantangan komputasi. Berdasarkan permasalahan tersebut, penelitian ini merumuskan masalah mengenai bagaimana mengimplementasikan dan mengevaluasi metode pembelajaran agar akurasi sistem ASR untuk bahasa Minangkabau dapat ditingkatkan. Tujuan dari penelitian ini adalah untuk mengadaptasi model bahasa universal ke dalam bahasa daerah secara efisien serta mengukur peningkatannya.

	Metodologi penelitian yang dilakukan berpusat pada proses \textit{fine-tuning} model \textit{pre-trained} OpenAI Whisper varian Base. Mengingat bahasa Minangkabau merupakan bahasa \textit{low-resource} dengan ketidakseimbangan frekuensi kemunculan fonem, penelitian ini menerapkan teknik \textit{Parameter-Efficient Fine-Tuning} (PEFT) melalui arsitektur \textit{Low-Rank Adaptation} (LoRA) untuk meminimalisasi beban komputasi selama pelatihan. Selain itu, digunakan pula optimasi berbasis \textit{Weighted Cross-Entropy Loss} guna menangani distribusi data yang tidak seimbang (\textit{data imbalance}) pada hierarki karakter, sehingga tebakan model tidak bias pada karakter yang umum saja. Evaluasi diukur secara kuantitatif dengan membandingkan nilai \textit{Word Error Rate} (WER) dan \textit{Character Error Rate} (CER) antara model \textit{zero-shot} bawaan model dasar dengan model hasil adaptasi.

	Hasil dan analisis observasi membuktikan bahwa proses \textit{fine-tuning} berhasil menekan persentase kesalahan prediksi secara substansial. Model Whisper yang telah di-\textit{fine-tuning} dengan modifikasi bobot mampu mengenali dan mentranskripsikan tuturan Minangkabau lebih akurat ketimbang generasi bawaannya. Di tingkat fonem, penerapan \textit{Weighted Cross-Entropy} membantu meluruskan kekeliruan pemetaan fonetik, meskipun sistem masih mendapati residu kesalahan ketika dihadapkan dengan penutur dari kawasan berdilek kental, seperti dari daerah Agam atau Lima Puluh Kota, yang memiliki serapan kata berbeda.

	Kesimpulan dari penelitian ini menegaskan bahwa pendekatan \textit{fine-tuning} menggunakan LoRA dipadukan dengan \textit{Weighted Cross-Entropy} pada OpenAI Whisper merupakan solusi yang sangat efektif bagi masalah ASR pada bahasa lokal. Perpaduan ini membuahkan performansi transkripsi yang jitu murni dengan mempertahankan efisiensi waktu komputasi. Sebagai saran untuk pengembangan selanjutnya, pengumpulan korpus data wicara yang lebih masif direkomendasikan sebagai prioritas utama. Perbanyakan data dapat dilakukan melintasi beragam kabupaten melalui metode \textit{crowdsourcing} penutur asli maupun ekstraksi arsip dari media siaran lokal agar model dapat terus berekspansi menangkap luasnya diksi verbal Minangkabau yang sesungguhnya.

\vfill
\clearpage